\documentclass[11pt,a4paper,twoside]{report}

% --- packages -----------------------------------------------------------------
\usepackage[utf8]{inputenc}
\usepackage[T1]{fontenc}
\usepackage{lmodern}
\usepackage[a4paper,margin=2.4cm,headheight=14pt]{geometry}
\usepackage{microtype}
\usepackage{xcolor}
\usepackage{listings}
\usepackage{hyperref}
\usepackage{enumitem}
\usepackage{tabularx}
\usepackage{booktabs}
\usepackage{longtable}
\usepackage{titlesec}
\usepackage{fancyhdr}
\usepackage{parskip}
\usepackage{amsmath,amssymb,amsthm}
\usepackage{multicol}
\usepackage{float}
\usepackage{caption}

% --- colours (Tokyo Night palette) --------------------------------------------
\definecolor{tnblue}{HTML}{7AA2F7}
\definecolor{tnred}{HTML}{F7768E}
\definecolor{tngreen}{HTML}{73DACA}
\definecolor{tncomment}{HTML}{565F89}
\definecolor{linkcolor}{HTML}{2C5AA0}
\definecolor{lightgray}{HTML}{EAEAEA}

% --- hyperref -----------------------------------------------------------------
\hypersetup{colorlinks=true,linkcolor=linkcolor,urlcolor=linkcolor,citecolor=linkcolor,
  pdftitle={sotFS System Manual},pdfauthor={Lucas Sotomayor},
  pdfsubject={Graph-Structured Filesystem with DPO Rewriting}}

% --- listings -----------------------------------------------------------------
\lstset{basicstyle=\ttfamily\footnotesize,numbers=left,numberstyle=\tiny\color{tncomment},
  numbersep=8pt,frame=single,framesep=4pt,rulecolor=\color{lightgray},
  backgroundcolor=\color{lightgray!30},showstringspaces=false,breaklines=true,
  tabsize=2,xleftmargin=1.5em,aboveskip=1em,belowskip=1em}

% --- titlesec -----------------------------------------------------------------
\titleformat{\chapter}[display]{\normalfont\huge\bfseries\color{tnblue}}
  {\chaptertitlename\ \thechapter}{20pt}{\Huge}
\titleformat{\section}{\normalfont\Large\bfseries\color{tnblue}}{\thesection}{1em}{}
\titleformat{\subsection}{\normalfont\large\bfseries\color{tnblue!80!black}}{\thesubsection}{1em}{}

% --- header / footer ----------------------------------------------------------
\pagestyle{fancy}\fancyhf{}
\fancyhead[LE]{\textit{sotFS System Manual}}
\fancyhead[RO]{\textit{\leftmark}}
\fancyfoot[C]{\thepage}
\renewcommand{\headrulewidth}{0.3pt}

% --- theorem environments -----------------------------------------------------
\theoremstyle{definition}
\newtheorem{definition}{Definition}[chapter]
\newtheorem{theorem}{Theorem}[chapter]
\newtheorem{invariant}{Invariant}[chapter]

% --- macros -------------------------------------------------------------------
\newcommand{\sotFS}{\texttt{sotFS}}
\newcommand{\sotX}{\texttt{sotX}}
\newcommand{\TG}{\mathrm{TG}}

% ==============================================================================
\begin{document}

\begin{titlepage}
\centering\vspace*{3cm}
{\Huge\bfseries\color{tnblue} sotFS System Manual\par}
\vspace{0.8cm}
{\Large Graph-Structured Filesystem with DPO Rewriting,\\
Bounded Treewidth, and Capability-Gated Projections\par}
\vspace{1.5cm}
{\large Version 1.0\par}
\vspace{0.5cm}
{\large Lucas Sotomayor\par}
\vspace{0.5cm}
{\large April 2026\par}
\vfill
{\small Comprehensive technical reference for the \sotFS{} filesystem.\\
Covers design, formal graph model, DPO rewriting rules, verification,\\
implementation, and architecture decisions.\par}
\end{titlepage}

\tableofcontents
\newpage

% ==============================================================================
\part{Design}
% ==============================================================================

\chapter{Introduction and Motivation}

Filesystem metadata is inherently a graph: inodes link to directories,
capabilities reference files, hard links create shared substructure, and
version histories form DAGs.  Yet every major filesystem---ext4, ZFS, Btrfs,
HAMMER2---treats metadata as flat tables or implicit trees, losing the graph
structure that could enable powerful queries, anomaly detection, and formal
reasoning about correctness.

\sotFS{} (Secure Object Transactional File System) makes this graph structure
explicit.  It represents filesystem metadata as a \textbf{typed graph} with six
node types and seven edge types, and defines POSIX operations as \textbf{Double
Pushout (DPO) graph rewriting rules}---algebraic transformations with
mathematically precise preconditions, postconditions, and gluing constraints.

\section{Contributions}
\begin{enumerate}
\item \textbf{Type Graph \TG{}:} A formal definition with six node types
  (Inode, Directory, Capability, Transaction, Version, Block) and seven
  edge types (contains, grants, delegates, derivesFrom, supersedes,
  pointsTo, hasXattr).
\item \textbf{DPO Rule Set:} Complete specification of POSIX operations as
  DPO graph rewriting rules with explicit gluing conditions.
\item \textbf{Treewidth Preservation:} Proof that DPO rules preserve bounded
  treewidth, enabling linear-time MSO queries via Courcelle's theorem.
\item \textbf{Curvature Monitor:} Incremental Ollivier--Ricci curvature
  computation for real-time structural anomaly detection.
\item \textbf{Deceptive Projections:} Capability-gated subgraph views as a
  graph endofunctor.
\item \textbf{GTXN Protocol:} Graph-level transaction manager composing SOT
  kernel transactions with graph-specific invariant checking.
\item \textbf{Coq Mechanization:} Machine-checked proofs that 6 DPO rules
  preserve all 5 structural invariants.
\end{enumerate}

\section{Comparison with Existing Systems}

\begin{center}\small
\begin{tabular}{lcccccc}
\toprule
\textbf{Feature} & \textbf{ext4} & \textbf{ZFS} & \textbf{Btrfs} & \textbf{HAMMER2} & \textbf{SquirrelFS} & \textbf{\sotFS{}} \\
\midrule
Metadata model & flat & Merkle & B-trees & radix & Rust types & \textbf{typed graph} \\
Formal ops     & none & none   & none    & none  & typestate  & \textbf{DPO rules} \\
Treewidth      & N/A  & N/A    & N/A     & N/A   & N/A        & \textbf{$\leq k$} \\
MSO queries    & N/A  & N/A    & N/A     & N/A   & N/A        & \textbf{O(n)} \\
Anomaly detect & none & none   & none    & none  & none       & \textbf{Ollivier--Ricci} \\
Crash consist. & journal & CoW & CoW+log & CoW   & typestate  & \textbf{DPO+WAL+typestate} \\
Verification   & tested & tested & tested & tested & typestate & \textbf{TLA+ $\to$ Coq} \\
\bottomrule
\end{tabular}
\end{center}

% ==============================================================================
\chapter{Architecture}

\section{sotFS Within the sotX Stack}

\begin{lstlisting}[language={},frame=single]
Application (POSIX binary via LUCAS/rump)
  | syscall (open, read, write, mkdir, ...)
LUCAS / rump VFS layer
  | Translates POSIX syscalls -> VfsOp IPC messages
sotFS Server Domain (Ring 3, userspace)
  |- VfsOp dispatcher
  |- GTXN engine (ADR-001)
  |   |- DPO rule matcher + gluing check
  |   |- Type graph (in-memory)
  |   |- Treewidth verifier
  |   +- Curvature monitor
  |- Persistence layer (WAL, bitmap, serialization)
  +- Bulk block-device capability (ADR-002)
sotX Microkernel (Ring 0): IPC | Scheduling | Frame alloc | Caps
Hardware (NVMe / virtio-blk)
\end{lstlisting}

\section{Data Flow: POSIX Call $\to$ Graph Rewrite $\to$ Persist}

Example: \texttt{mkdir("/data/newdir", 0o755)}:
\begin{enumerate}
\item LUCAS receives mkdir syscall, sends \texttt{VfsOp::Mkdir} IPC to \sotFS{}.
\item Dispatcher invokes DPO rule MKDIR.
\item GTXN begins: match Directory for ``/data'', check gluing conditions,
  begin SOT Tier~1 transaction.
\item Apply MKDIR rule: add Inode, Directory, three contains edges
  (entry, ``.'', ``..''). Verify link\_count consistency and treewidth bound.
\item GTXN commits: WAL write (atomic), bitmap update, curvature update,
  provenance emit.
\item Reply to LUCAS with success and new inode ID.
\end{enumerate}

% ==============================================================================
\chapter{Type Graph Formalization}

\begin{definition}[Type Graph]
The \sotFS{} type graph is a tuple
$\TG = (V, E, \mathrm{src}, \mathrm{tgt}, \tau_V, \tau_E, \mathrm{attr})$
where $V$ is a finite set of nodes, $E$ a finite set of edges,
$\mathrm{src}, \mathrm{tgt}: E \to V$ map edges to endpoints,
$\tau_V: V \to \mathrm{NodeType}$ and $\tau_E: E \to \mathrm{EdgeType}$
assign types, and $\mathrm{attr}: V \cup E \to \mathrm{Attributes}$
assigns typed attributes.

\medskip\noindent
$\mathrm{NodeType} = \{\text{Inode, Directory, Capability, Transaction, Version, Block}\}$\\
$\mathrm{EdgeType} = \{\text{contains, grants, delegates, derivesFrom, supersedes, pointsTo, hasXattr}\}$
\end{definition}

\section{Node Types}

\subsection{Inode}
Attributes: \texttt{id:~OID}, \texttt{vtype}, \texttt{perms}, \texttt{uid}, \texttt{gid},
\texttt{size}, \texttt{link\_count}, \texttt{ctime}, \texttt{mtime}, \texttt{atime}.

\textbf{Invariant I\textsubscript{2}:} $\texttt{link\_count} = |\{e \in E : \tau_E(e)=\text{contains} \wedge \mathrm{tgt}(e)=i \wedge e.\text{name} \neq \text{``..''}\}|$

\subsection{Directory}
Attributes: \texttt{id:~OID}. Paired 1:1 with an Inode of $\texttt{vtype}=\text{Directory}$ via a ``\texttt{.}'' edge.

\textbf{Invariant I\textsubscript{3}:} Every Directory $d$ has exactly one ``\texttt{.}'' contains edge to its paired Inode.

\textbf{Invariant I\textsubscript{4}:} Names on outgoing contains edges from $d$ are unique (set, not multiset).

\subsection{Capability}
Attributes: \texttt{cap\_id}, \texttt{rights} (bitmask: R/W/X/Grant/Revoke), \texttt{epoch}.

\textbf{Invariant I\textsubscript{5}:} Exactly one outgoing \texttt{grants} edge, at most one incoming \texttt{delegates} edge.

\subsection{Transaction, Version, Block}
Transaction: models GTXN state machine (\texttt{Active/Preparing/Committed/Aborted}).
Version: snapshot lineage node. Block: contiguous sector extent with \texttt{refcount}.

\section{Edge Types and Constraints}

\begin{center}\small
\begin{tabular}{llll}
\toprule
\textbf{Edge} & \textbf{Source $\to$ Target} & \textbf{Attributes} & \textbf{Key Constraint} \\
\midrule
contains & Directory $\to$ Inode & name & C\textsubscript{1}: unique names per dir \\
grants & Capability $\to$ Inode & rights & C\textsubscript{2}: $\text{rights} \subseteq \text{cap.rights}$ \\
delegates & Cap $\to$ Cap & -- & C\textsubscript{3}: monotonic attenuation; C\textsubscript{4}: forest \\
derivesFrom & Version $\to$ Version & -- & C\textsubscript{5}: DAG \\
supersedes & Inode $\to$ Inode & -- & C\textsubscript{6}: acyclic \\
pointsTo & Inode $\to$ Block & offset & C\textsubscript{7}: non-overlapping extents \\
hasXattr & Inode $\to$ XAttr & -- & 1:1 per (inode, namespace, name) \\
\bottomrule
\end{tabular}
\end{center}

\section{Global Invariants}

\begin{invariant}[G\textsubscript{1} -- No orphan inodes]
$\forall i \in V, \tau_V(i) = \text{Inode} \implies \text{reachable}(\text{root}, i) \vee i.\text{link\_count} = 0$
\end{invariant}

\begin{invariant}[G\textsubscript{2} -- No dangling edges]
$\forall e \in E: \mathrm{src}(e) \in V \wedge \mathrm{tgt}(e) \in V$
\end{invariant}

\begin{invariant}[G\textsubscript{3} -- Link count consistency]
$\forall i: i.\text{link\_count} = |\{e : \tau_E(e)=\text{contains} \wedge \mathrm{tgt}(e)=i \wedge e.\text{name} \neq \text{``..''}\}|$
\end{invariant}

\begin{invariant}[G\textsubscript{4} -- Capability monotonicity]
Along any delegates path $c_0 \to c_1 \to \cdots \to c_n$:
$c_n.\text{rights} \subseteq \cdots \subseteq c_0.\text{rights}$
\end{invariant}

\begin{invariant}[G\textsubscript{5} -- No directory cycles]
The contains subgraph restricted to Directory$\to$Inode(vtype=Directory) links (excluding ``\texttt{.}'' and ``\texttt{..}'') is acyclic.
\end{invariant}

\begin{invariant}[G\textsubscript{7} -- Treewidth bound]
$\mathrm{tw}(\TG) \leq k$ for configurable $k$ (see Theorem~\ref{thm:tw}).
\end{invariant}

% ==============================================================================
\chapter{DPO Rules for POSIX Operations}

A DPO rule $\rho = (L \xleftarrow{l} K \xrightarrow{r} R)$ transforms host graph $G$ via match $m: L \to G$ into result $H$, subject to gluing conditions (identification and dangling).

\section{Mutating Rules}

\subsection{CREATE(dir, name, vtype) $\to$ new\_oid}
$L = \{d\text{:Dir}\}$, $K = \{d\}$, $R = \{d, i\text{:Inode[lc=1]}, e\text{:contains}(d,i)\text{[name]}\}$.
If vtype=Directory, additionally creates Directory node, ``\texttt{.}'' and ``\texttt{..}'' edges.

\textbf{Gluing:} GC-CREATE-1 (no duplicate name), GC-CREATE-2 (dir exists).

\subsection{WRITE(inode, offset, data)}
Adds/replaces Block + pointsTo edge. CoW check on shared blocks.

\subsection{MKDIR(dir, name)}
Special CREATE adding Inode(lc=2) + Directory + 3 contains edges.

\subsection{RMDIR(dir, name)}
Removes Directory + Inode + 3 edges. \textbf{Gluing:} GC-RMDIR-1 (directory empty).

\subsection{LINK(dir, name, target)}
Adds contains edge, increments link\_count. \textbf{Gluing:} GC-LINK-2 (target is Regular), GC-LINK-3 ($\text{lc} < \text{LINK\_MAX}$).

\subsection{UNLINK(dir, name)}
Removes contains edge, decrements link\_count. GC if lc reaches 0.

\subsection{RENAME(src\_dir, src\_name, dst\_dir, dst\_name)}
Four cases: same/cross dir $\times$ target exists/absent.  Cross-dir requires SOT Tier~2 (2PC).
\textbf{Gluing:} GC-RENAME-2 (cycle prevention via \texttt{is\_ancestor} check).

\section{Treewidth Impact per Rule}
Only LINK can increase treewidth, bounded by GC-LINK-3. All other rules are tw-neutral or tw-decreasing.

% ==============================================================================
\part{Theory}
% ==============================================================================

\chapter{Bounded Treewidth and MSO Queries}

\begin{theorem}[Treewidth Bound]\label{thm:tw}
Let $\TG$ be produced by any finite DPO rule sequence from CREATE-ROOT.
If $\text{link\_count} \leq \text{LINK\_MAX}$ per inode, then
$\mathrm{tw}(\TG) \leq \text{LINK\_MAX} + O(1)$.
\end{theorem}

\textbf{Proof sketch.} The contains subgraph without hard links is a forest (tw=1). Each hard link adds one back-edge, increasing tw by at most 1. The grants, delegates, derivesFrom, and pointsTo subgraphs each contribute tw$\leq$1. Combined: $\mathrm{tw}(\TG) \leq \text{LINK\_MAX} + 5$.

\textbf{Corollary.} Without hard links: $\mathrm{tw}(\TG) \leq 6$.

\section{Courcelle's Theorem Application}
Since $\mathrm{tw}(\TG)$ is bounded, every MSO\textsubscript{2}-definable query runs in $O(n)$ time.
Examples: capability reachability, cycle detection, orphan detection, version ancestry.

% ==============================================================================
\chapter{Ollivier--Ricci Curvature Monitor}

For edge $(u,v)$: $\kappa(u,v) = 1 - W_1(\mu_u, \mu_v) / d(u,v)$ where $\mu_x$ is the lazy random walk measure with parameter $\alpha = 0.5$.

\textbf{Normal patterns:} CREATE decreases parent curvature slightly; UNLINK increases it; WRITE has no effect; steady-state $\kappa \approx -0.5$ to $0$.

\textbf{Anomalous signatures:} Mass creation $\to$ positive curvature spike; symlink bomb $\to$ strong positive on target; deep chain $\to$ $\kappa \approx -1.0$.

\textbf{Incremental:} After GTXN commit, recompute only the 2-hop neighborhood of changed edges. Amortized $O(d_{\max}^2)$ per rule application.

% ==============================================================================
\chapter{Deceptive Subgraph Views}

A projection $\pi$ for domain $D$ is a typed graph morphism
$\pi: \TG|_D \to \TG_D$ that may omit, add, or relabel nodes/edges.
The deception functor $D_{FS}: \textbf{TGraph} \to \textbf{TGraph}$ extends
the SOT deception model to graph structure.

Consistency requirement: $D_{FS}(\TG)$ must satisfy all global invariants G\textsubscript{1}--G\textsubscript{7}.
Fabricated content is checkable via MSO queries in linear time.

% ==============================================================================
\chapter{Capabilities and Transactions}

\section{Capabilities as Graph Edges}
DERIVE: creates child capability with $\text{rights} \subseteq \text{parent.rights}$ via \texttt{delegates} edge.
REVOKE: epoch-based, removes capability node and all delegates descendants.

\section{Graph Transactions (GTXN)}
\textbf{Protocol:} begin $\to$ verify\_gluing $\to$ verify\_treewidth $\to$ SOT::begin $\to$ apply\_rule $\to$ verify\_postconditions $\to$ commit/rollback.

\textbf{Tier mapping:} Most ops use SOT Tier~1. Cross-directory RENAME uses Tier~2 (2PC).

% ==============================================================================
\part{Verification}
% ==============================================================================

\chapter{Multi-Layer Verification}

\section{Layer 1: TLA+ Model Checking (Bounded)}
4 specifications verified: \texttt{sotfs\_graph}, \texttt{sotfs\_transactions}, \texttt{sotfs\_capabilities}, \texttt{sotfs\_crash}. 20+ invariants model-checked across small/medium/large configurations.

\section{Layer 2: Crash Refinement (Tested)}
6 integration tests with redb ACID backend verify round-trip persistence, atomic overwrite, and complex operation survival.

\section{Layer 3: Rust Typestate (Compile-Time)}
\texttt{InodeHandle<State>}, \texttt{TxHandle<State>}, \texttt{DirHandle<State>}, \texttt{CapHandle} enforce state machine transitions at compile time. \texttt{\#[must\_use]} prevents silent transaction drop.

\section{Layer 4: Loom Concurrency (Bounded)}
3 exhaustive interleaving tests: transaction isolation, rollback correctness, conflicting write serialization.

\section{Layer 5: Coq Mechanized Proofs (Unbounded)}

\begin{center}\small
\begin{tabular}{llll}
\toprule
\textbf{Module} & \textbf{DPO Rule} & \textbf{Invariants} & \textbf{Status} \\
\midrule
\texttt{SotfsGraph.v} & (definitions) & init\_graph\_well\_formed & Complete \\
\texttt{DpoCreate.v} & create\_file & all 5 & Complete (0 admits) \\
\texttt{DpoUnlink.v} & unlink\_keep & all 5 & Complete (0 admits) \\
\texttt{DpoRename.v} & rename\_same\_dir & all 5 & Complete (0 admits) \\
\texttt{DpoMkdir.v} & mkdir & all 5 & Complete (1 side cond.) \\
\texttt{DpoLink.v} & hard\_link & all 5 & Complete (0 admits) \\
\texttt{DpoRmdir.v} & rmdir & 3 of 5 & Partial \\
\bottomrule
\end{tabular}
\end{center}

\textbf{Key lemma:} \texttt{count\_remove\_matching} establishes that removing one predicate-satisfying element decreases \texttt{count\_occ\_pred} by exactly one, closing LinkCountConsistent for Unlink and Rename.

\section{Assumptions (Not Verified)}
\textbf{Hardware:} CPU correctness (x86\_64 TSO), storage durability (fsync honored), no bit rot.
\textbf{Compiler:} Rust compiler correctness, no UB in safe code, optimization preserves semantics.
\textbf{Design:} IPC reliability (verified by \texttt{ipc.tla}), capability unforgery, single-writer per GTXN.

% ==============================================================================
\chapter{TLC Model Checking Bounds}

Three configuration tiers per spec.  The runner script \texttt{run\_tlc.sh} drives TLC across all combinations.

\begin{center}\small
\begin{tabular}{llccc}
\toprule
\textbf{Spec} & \textbf{Key Constant} & \textbf{Small} & \textbf{Medium} & \textbf{Large} \\
\midrule
sotfs\_graph & InodeIds & 5 & 6 & 8 \\
             & DirIds   & 3 & 4 & 5 \\
             & MaxOps   & 6 & 8 & 10 \\
\midrule
sotfs\_transactions & GTxnIds & 2 & 3 & 4 \\
                    & SOIds   & 3 & 4 & 5 \\
                    & MaxOps  & 8 & 10 & 12 \\
\midrule
sotfs\_capabilities & CapIds    & 4 & 5 & 5 \\
                    & InodeIds  & 2 & 3 & 3 \\
                    & MaxOps    & 6 & 8 & 8 \\
\midrule
sotfs\_crash & SOIds  & 2 & 3 & 4 \\
             & MaxOps & 8 & 10 & 12 \\
\bottomrule
\end{tabular}
\end{center}

% ==============================================================================
\part{Implementation}
% ==============================================================================

\chapter{Crate Structure}

\begin{center}\small
\begin{tabular}{lll}
\toprule
\textbf{Crate} & \textbf{Description} & \textbf{Key Types} \\
\midrule
\texttt{sotfs-graph}   & Type graph data structures & \texttt{TypeGraph}, \texttt{Arena<T,N>}, \texttt{RcuGraph} \\
\texttt{sotfs-ops}     & DPO rewriting rules        & \texttt{create\_file}, \texttt{mkdir}, \texttt{symlink}, \texttt{fsck} \\
\texttt{sotfs-storage} & Persistence via redb        & \texttt{save()}, \texttt{load()} \\
\texttt{sotfs-tx}      & GTXN manager               & \texttt{with\_transaction}, commit/rollback \\
\texttt{sotfs-monitor} & Anomaly monitors            & treewidth, curvature, deception \\
\texttt{sotfs-fuse}    & FUSE bridge                 & Full POSIX VFS \\
\bottomrule
\end{tabular}
\end{center}

\section{Arena Allocator}
Fixed-capacity heap-backed pool (\texttt{Box<[MaybeUninit<T>]>}) with O(1) alloc/dealloc/lookup.
Free-list slot reuse prevents fragmentation.  65K node slots, 131K edge slots.
TypeGraph struct: $\sim$744 bytes (pointers + counters).

\section{Epoch-Based RCU}
Lock-free reads via two TypeGraph copies with atomic swap.  8 concurrent reader slots.
Writers: clone active $\to$ inactive, mutate, swap, wait for epoch drain.

\section{DynamicTreewidth}
Incremental treewidth via cached elimination ordering.
$<$20\% dirty nodes: partial re-elimination.  $\geq$20\%: full greedy min-degree.

% ==============================================================================
\chapter{New Features}

\section{Extended Attributes (xattrs)}
XAttr nodes connected via HasXattr edges.  Operations: setxattr, getxattr, removexattr, listxattr.
Four namespaces (User, System, Security, Trusted).  64KB max value.

\section{Symlinks}
\texttt{VnodeType::Symlink} with target in \texttt{symlink\_targets: BTreeMap}.
DPO rule identical to CREATE except vtype=Symlink, permissions=0777.

\section{POSIX ACLs $\leftrightarrow$ Capability Graph}
ACL entries map to Grants edges: \texttt{ACL\_USER\_OBJ} $\to$ \texttt{Grants(cap\_owner, inode, rights)}.
Default: 3-entry ACL synthesized from permission bits.

\section{Hierarchical Quotas}
Per-subtree \texttt{Quota} with inode/byte limits.  Enforcement walks ``\texttt{..}'' edges upward.
Propagation: $O(\text{depth})$ amortized.

\section{fsck Verifier}
Checks all 5 structural invariants + block refcount + orphan detection + xattr integrity.
Returns \texttt{FsckReport} with typed errors and warnings.

% ==============================================================================
\chapter{Test Summary}

\begin{center}\small
\begin{tabular}{lrl}
\toprule
\textbf{Category} & \textbf{Count} & \textbf{Method} \\
\midrule
Arena allocator     & 16 & Unit tests \\
RCU concurrency     &  9 & Multi-thread \\
Graph invariants    &  3 & Unit tests \\
Typestate           &  6 & Unit + compile-fail \\
DPO operations      & 35 & Unit tests \\
Property-based      & 10 & proptest (10K iter) \\
Storage/crash       &  7 & Integration (redb) \\
Transaction         &  3 & Unit tests \\
Loom concurrency    &  3 & Exhaustive \\
Treewidth           & 20 & Unit + adversarial \\
Curvature           & 14 & Unit + incremental \\
Deception           &  5 & Projections \\
Adversarial         & 13 & Security scenarios \\
xattrs/symlinks/ACL & 10 & Unit tests \\
fsck                &  2 & Unit tests \\
\midrule
\textbf{Total}      & \textbf{147} & \textbf{All passing} \\
\bottomrule
\end{tabular}
\end{center}

% ==============================================================================
\part{Architecture Decisions}
% ==============================================================================

\chapter{ADR-001: Graph-Level Transaction Manager}
\textbf{Status:} Accepted (2026-04-10). \textbf{Decision:} Separate userspace GTXN wrapping SOT Tier~1/2 transactions.  Preserves kernel simplicity; checks graph invariants in userspace.  Rejected alternatives: direct SOT use (no invariant checking), unified kernel manager (violates microkernel principle).

\chapter{ADR-002: Block Capability Interpretation}
\textbf{Status:} Accepted. \textbf{Decision:} sotFS server holds bulk block-device capability; clients get attenuated inode-SO caps.  Matches existing VFS pattern.  Rejected: per-inode block caps (capability explosion).

\chapter{ADR-003: Bootstrap Procedure}
\textbf{Status:} Accepted. \textbf{Decision:} \texttt{sotfs\_format()} is a DPO rule application (CREATE-ROOT).  Initializes empty graph with root inode, directory, capability, superblock.

\chapter{ADR-004: Rump Kernel Interaction}
\textbf{Status:} Accepted. \textbf{Decision:} \sotFS{} registers as \texttt{FsType::SotFs} in rump VFS mount table, coexisting with FFS.

% ==============================================================================
\appendix
\chapter{VfsOp $\to$ DPO Rule Mapping}

\begin{center}\small
\begin{tabular}{llll}
\toprule
\textbf{VfsOp} & \textbf{DPO Rule} & \textbf{Tier} & \textbf{Graph Mutation} \\
\midrule
Open    & OPEN (query)  & 0   & None \\
Close   & CLOSE (query) & 0   & May trigger GC \\
Read    & READ (query)  & 0   & None \\
Write   & WRITE         & 1   & Add/replace Block \\
Mkdir   & MKDIR         & 1   & Add Inode + Dir + edges \\
Rmdir   & RMDIR         & 1   & Remove Inode + Dir + edges \\
Unlink  & UNLINK        & 1   & Remove edge, maybe Inode \\
Rename  & RENAME        & 1/2 & Move edge \\
Stat    & STAT (query)  & 0   & None \\
Readdir & READDIR       & 0   & Enumerate contains edges \\
Fsync   & ---           & --- & Force WAL flush \\
\bottomrule
\end{tabular}
\end{center}

\chapter{Category Theory Integration}

\sotFS{} introduces category \textbf{TGraph} (typed graphs, typed graph morphisms) alongside the existing SOT categories.

\textbf{Functors:}
\begin{itemize}
\item $G: \mathbf{SO} \to \mathbf{TGraph}$ --- maps Secure Objects to graph nodes (faithful).
\item $Rw: \mathbf{Tx} \to \mathbf{TGraph}$ --- maps transactions to graph transformations (composition-preserving).
\item $D_{FS}: \mathbf{TGraph} \to \mathbf{TGraph}$ --- deception endofunctor (invariant-preserving).
\end{itemize}

\textbf{Natural transformations:}
$\eta: \mathrm{Id} \Rightarrow \mathrm{WAL} \circ \mathrm{Id}$ (persistence),
$\pi_D: \mathrm{Id} \Rightarrow D_{FS}$ (projection).

\end{document}
